\documentclass[12pt]{article}
\usepackage[a4paper,margin=1in]{geometry}
\usepackage[utf8]{inputenc}
\usepackage[T1]{fontenc}
\usepackage{amsmath,amssymb}
\usepackage{textcomp}
\usepackage{parskip}
\begin{document}
\begin{center}
{\Large \textbf{2019 STEP II}}
\end{center}

\subsection*{Page 1}
\begin{flushleft}
\bigskip\textbf{Section A: Pure Mathematics}\par
\bigskip\textbf{Question 1.} Let f(x) = (x p)g(x), where g is a polynomial. Show that the tangent to the curve y = f(x)\par
\[- \leq\]
at the point with x = a, where a = p, passes through the point (p, 0) if and only if g (a) = 0.\par
\[\leq\]
The curve C has equation\par
\[y = A(x p)(x q)(x r),\]
- - -\par
\[where p, q and r are constants with p < q < r, and A is a non-zero constant.\]
\textbf{(i)} The tangent to C at the point with x = a, where a = p, passes through the point (p, 0).\par
\[\leq\]
\[Show that 2a = q + r and find an expression for the gradient of this tangent in terms\]
of A, q and r.\par
\textbf{(ii)} The tangent to C at the point with x = c, where c = r, passes through the point (r, 0).\par
\[\leq\]
Show that this tangent is parallel to the tangent in part (i) if and only if the tangent\par
\[to C at the point with x = q does not meet the curve again.\]
\[\leq\]
\bigskip\textbf{Question 2.} The function f satisfies f(0) = 0 and f (t) > 0 for t > 0 . Show by means of a sketch that,\par
\[for x > 0,\]
x f(x)\par
f(t) dt + f\par
-1(y)\par
\[dy = xf(x).\]
\[\leq \leq\]
0 0\par
\textbf{(i)} The (real) function g is defined, for all t, by\par
\[g(t) + g(t) = t.\]
\[\geq (\]
\[\leq\]
\[Prove that g(0) = 0 , and that g (t) > 0 for all t.\]
Evaluate g(t) dt.\par
\[\leq\]
\textbf{(ii)} The (real) function h is defined, for all t, by\par
\[h(t) + h(t) = t + 2.\]
\[\geq (\]
Evaluate h(t) dt .\par
\[\leq\]
\end{flushleft}
\subsection*{Page 2}
\begin{flushleft}
\bigskip\textbf{Question 3.} For any two real numbers x and x , show that\par
1 2\par
\[x + x \leq x + x .\]
1 2 1 2\par
| | | | | |\par
Show further that, for any real numbers x , x , . . . , x ,\par
1 2 n\par
\[x + x + + x \leq x + x + + x .\]
1 2 n 1 2 n\par
| \cdots | | | | | \cdots | |\par
\textbf{(i)} The polynomial f is defined by\par
\[f(x) = 1 + a 1 x + a 2 x^2 + + a n-1 xn-1 + xn\]
\cdots\par
where the coefficients are real and satisfy a \leq A for i = 1, 2, . . . , n 1, where A \geq 1 .\par
i\par
| | -\par
\[(a) If x < 1, show that\]
| |\par
A x\par
\[f(x) 1 \leq | | .\]
| - | 1 x\par
- | |\par
\[(b) Let ω be a real root of f, so that f(ω) = 0 . In the case ω < 1, show that\]
| |\par
\[\leq ω \leq 1 + A. ( )\]
1 + A | | ∗\par
\[(c) Show further that the inequalities ( ) also hold if ω \geq 1.\]
∗ | |\par
\textbf{(ii)} Find the integer root or roots of the quintic equation\par
\[135x5 135x4 100x3 91x2 126x + 135 = 0.\]
- - - -\par
\end{flushleft}
\subsection*{Page 3}
\begin{flushleft}
\bigskip\textbf{Question 4.} You are not required to consider issues of convergence in this question.\par
n\par
For any sequence of numbers a , a , . . . , a , . . . , a , the notation a denotes the\par
1 2 m n i\par
\[i)=m\]
product a a a .\par
m m+1 n\par
\cdots\par
\textbf{(i)} Use the identity 2 cos x sin x = sin(2x) to evaluate the product cos( \pi ) cos( 2\pi ) cos( 4\pi ).\par
9 9 9\par
\textbf{(ii)} Simplify the expression\par
n\par
x\par
\[cos (0 < x < 1 \pi).\]
2k 2\par
\[k)=0 \int \sum\]
\[Using differentiation, or otherwise, show that, for 0 < x < 1 \pi,\]
n\par
1 x 1 x\par
\[tan = cot 2 cot(2x).\]
2k 2k 2n 2n -\par
\[(k=0 \int \sum \int \sum\]
sin θ tan θ\par
\textbf{(iii)} Using the results lim = 1 and lim = 1 , show that\par
\[θ\to0 θ θ\to0 θ\]
\infty\par
x sin x\par
\[cos =\]
2k x\par
\[k)=1 \int \sum\]
and evaluate\par
\infty\par
\[1 \pi\]
tan .\par
2j-2 2j\par
\[(j=2 \int \sum\]
\end{flushleft}
\subsection*{Page 4}
\begin{flushleft}
\bigskip\textbf{Question 5.} The sequence u , u , . . . is said to be a constant sequence if u = u for n = 0, 1, 2, . . . .\par
0 1 n n+1\par
\[The sequence is said to be a sequence of period 2 if u = u for n = 0, 1, 2, . . . and the\]
n n+2\par
sequence is not constant.\par
\textbf{(i)} A sequence of real numbers is defined by u = a and u = f(u ) for n = 0, 1, 2, . . . ,\par
0 n+1 n\par
where\par
\[f(x) = p + (x p)x,\]
-\par
and p is a given real number.\par
Find the values of a for which the sequence is constant.\par
\[Show that the sequence has period 2 for some value of a if and only if p > 3 or p < 1 .\]
-\par
\textbf{(ii)} A sequence of real numbers is defined by u = a and u = f(u ) for n = 0, 1, 2, . . . ,\par
0 n+1 n\par
where\par
\[f(x) = q + (x p)x,\]
-\par
and p and q are given real numbers.\par
Show that there is no value of a for which the sequence is constant if and only if\par
\[f(x) > x for all x.\]
Deduce that, if there is no value of a for which the sequence is constant, then there is\par
no value of a for which the sequence has period 2.\par
Is it true that, if there is no value of a for which the sequence has period 2, then there\par
is no value of a for which the sequence is constant?\par
\end{flushleft}
\subsection*{Page 5}
\begin{flushleft}
6 Note: You may assume that if the functions y (x) and y (x) both satisfy one of the differ-\par
1 2\par
ential equations in this question, then the curves y = y (x) and y = y (x) do not intersect.\par
1 2\par
\textbf{(i)} Find the solution of the differential equation\par
dy\par
\[= y + x + 1\]
dx\par
\[that has the form y = mx + c, where m and c are constants.\]
\[Let y (x) be the solution of this differential equation with y (0) = k. Show that any\]
3 3\par
\[stationary point on the curve y = y (x) lies on the line y = x 1. Deduce that\]
- -\par
\[solution curves with k < 2 cannot have any stationary points.\]
-\par
Show further that any stationary point on the solution curve is a local minimum.\par
\[Use the substitution Y = y + x to solve the differential equation, and sketch, on the\]
\[same axes, the solutions with k = 0, k = 2 and k = 3.\]
- -\par
\textbf{(ii)} Find the two solutions of the differential equation\par
dy\par
\[= x^2 + y^2 2xy 4x + 4y + 3\]
dx - -\par
\[that have the form y = mx + c.\]
\[Let y (x) be the solution of this differential equation with y (0) = 2. (Do not attempt\]
4 4\par
-\par
to find this solution.)\par
\[Show that any stationary point on the curve y = y (x) lies on one of two lines that you\]
should identify. What can be said about the gradient of the curve at points between\par
these lines?\par
\[Sketch the curve y = y (x). You should include on your sketch the two straight line\]
solutions and the two lines of stationary points.\par
\end{flushleft}
\subsection*{Page 6}
\begin{flushleft}
\bigskip\textbf{Question 7.} (i) The points A, B and C have position vectors a, b and c, respectively. Each of these\par
\[vectors is a unit vector (so a . a = 1 , for example) and\]
\[a + b + c = 0.\]
\[Show that a . b = 1 . What can be said about the triangle ABC? You should justify\]
- 2\par
your answer.\par
\textbf{(ii)} The four distinct points A (i = 1, 2, 3, 4) have unit position vectors a and\par
i i\par
\[a = 0.\]
i\par
\[(i=1\]
\[Show that a . a = a . a .\]
1 2 3 4\par
(a) Given that the four points lie in a plane, determine the shape of the quadrilateral\par
with vertices A , A , A and A .\par
1 2 3 4\par
(b) Given instead that the four points are the vertices of a regular tetrahedron, find\par
the length of the sides of this tetrahedron.\par
\bigskip\textbf{Question 8.} The domain of the function f is the set of all 2 2 matrices and its range is the set of real\par
\[\times\]
numbers. Thus, if M is a 2 2 matrix, then f(M) R.\par
\[\times \in\]
\[The function f has the property that f(MN) = f(M)f(N) for any 2 2 matrices M and N.\]
\[\times\]
\textbf{(i)} You are given that there is a matrix M such that f(M) = 0. Let I be the 2 2 identity\par
\[\leq \times\]
\[matrix. By considering f(MI), show that f(I) = 1.\]
0 1\par
\textbf{(ii)} Let J = . You are given that f(J) = 1. By considering J2, evaluate f(J).\par
\[) 1 0 [ \leq\]
Using J, show that, for any real numbers a, b, c and d,\par
a b c d d c\par
\[f = f = f .\]
) c d [ - ) a b [ ) b a [\par
\[\int \sum \int \sum \int \sum\]
1 0\par
\textbf{(iii)} Let K = where k R. Use K to show that, if the second row of the\par
) 0 k [ \in\par
\[matrix A is a multiple of the first row, then f(A) = 0.\]
1 1\par
\textbf{(iv)} Let P = . By considering the matrices P2, P -1, and K -1PK for suitable\par
) 0 1 [\par
values of k, evaluate f(P).\par
\end{flushleft}
\subsection*{Page 7}
\begin{flushleft}
\bigskip\textbf{Section B: Mechanics}\par
\bigskip\textbf{Question 9.} A particle P is projected from a point O on horizontal ground with speed u and angle of\par
\[projection α, where 0 < α < 1 \pi.\]
\[2\sqrt\{2\}\]
\textbf{(i)} Show that if sin α < , then the distance OP is increasing throughout the flight.\par
\[2\sqrt\{2\}\]
\[Show also that if sin α > , then OP will be decreasing at some time before the\]
particle lands.\par
\textbf{(ii)} At the same time as P is projected, a particle Q is projected horizontally from O\par
with speed v along the ground in the opposite direction from the trajectory of P . The\par
ground is smooth. Show that if\par
\[2\sqrt\{2v\} > (sin α 2\sqrt\{2\} cos α)u,\]
-\par
then QP is increasing throughout the flight of P .\par
\bigskip\textbf{Question 10.} A small light ring is attached to the end A of a uniform rod AB of weight W and length 2a.\par
The ring can slide on a rough horizontal rail.\par
One end of a light inextensible string of length 2a is attached to the rod at B and the other\par
end is attached to a point C on the rail so that the rod makes an angle of θ with the rail,\par
\textasciicircum{}\circ\par
\[where 0 < θ < 90 . The rod hangs in the same vertical plane as the rail.\]
A force of kW acts vertically downwards on the rod at B and the rod is in equilibrium.\par
\textbf{(i)} You are given that the string will break if the tension T is greater than W . Show that\par
(assuming that the ring does not slip) the string will break if\par
\[2k + 1 > 4 sin θ.\]
\textbf{(ii)} Show that (assuming that the string does not break) the ring will slip if\par
\[2k + 1 > (2k + 3)μ tan θ,\]
where μ is the coefficient of friction between the rail and the ring.\par
\textbf{(iii)} You are now given that μ tan θ < 1.\par
Show that, when k is increased gradually from zero, the ring will slip before the string\par
breaks if\par
2 cos θ\par
\[μ < .\]
1 + 2 sin θ\par
\end{flushleft}
\subsection*{Page 8}
\begin{flushleft}
\bigskip\textbf{Section C: Probability and Statistics}\par
\bigskip\textbf{Question 11.} (i) The three integers n , n and n satisfy 0 < n < n < n and n + n > n . Find\par
1 2 3 1 2 3 1 2 3\par
\[the number of ways of choosing the pair of numbers n and n in the cases n = 9 and\]
1 2 3\par
\[n = 10.\]
\[Given that n = 2n + 1, where n is a positive integer, write down an expression (which\]
you need not prove is correct) for the number of ways of choosing the pair of numbers\par
n and n . Simplify your expression.\par
1 2\par
\[Write down and simplify the corresponding expression when n = 2n , where n is a\]
positive integer.\par
\textbf{(ii)} You have N rods, of lengths 1, 2, 3, . . . , N (one rod of each length). You take the\par
rod of length N , and choose two more rods at random from the remainder, each choice\par
\[of two being equally likely. Show that, in the case N = 2n + 1 where n is a positive\]
integer, the probability that these three rods can form a triangle (of non-zero area) is\par
n 1\par
- .\par
2n 1\par
-\par
\[Find the corresponding probability in the case N = 2n, where n is a positive integer.\]
\textbf{(iii)} You have 2M +1 rods, of lengths 1, 2, 3, . . . , 2M +1 (one rod of each length), where M\par
is a positive integer. You choose three at random, each choice of three being equally\par
likely. Show that the probability that the rods can form a triangle (of non-zero area)\par
is\par
(4M + 1)(M 1)\par
- .\par
2(2M + 1)(2M 1)\par
-\par
K\par
\[Note: k^2 = 1 K(K + 1)(2K + 1).\]
\[k]=1\]
\end{flushleft}
\subsection*{Page 9}
\begin{flushleft}
\bigskip\textbf{Question 12.} The random variable X has the probability density function on the interval [0, 1]:\par
\[nxn-1 0 \leq x \leq 1,\]
\[f(x) = ⎧\]
⎨0 elsewhere,\par
⎩\par
where n is an integer greater than 1.\par
\textbf{(i)} Let μ = E(X). Find an expression for μ in terms of n, and show that the variance, σ2,\par
of X is given by\par
n\par
\[σ2 = .\]
(n + 1)2(n + 2)\par
\textbf{(ii)} In the case n = 2, show without using decimal approximations that the interquartile\par
range is less than 2σ.\par
\textbf{(iii)} Write down the first three terms and the (k + 1)th term (where 0 \leq k \leq n) of the\par
binomial expansion of (1 + x) n in ascending powers of x.\par
\[By setting x = , show that μ is less than the median and greater than the lower\]
n\par
quartile.\par
Note: You may assume that\par
1 1 1\par
\[1 + + + + < 4.\]
1! 2! 3! \cdots\par
\end{flushleft}
\end{document}
